% 题干转写：完成
% 解析转写：完成
% 原题与答案均无独立图示；下图为解析辅助示意
\section{第39届复赛第2题：通电螺线管弹簧的平衡与等效劲度}

\subsection{题干}
某长直弹簧由涂绝缘漆的磷铜细丝绕成 $N$ 匝，可视为横截面半径为 $r$ 的均匀密绕长直螺线管。弹簧原长为 $l_0\gg r$，劲度系数为 $k$；形变中横截面半径不变，忽略边缘效应、漏磁和重力，真空磁导率为 $\mu_0$。

（1）用恒流源使电流保持为 $I_0$，起初以外力使弹簧不变形，缓慢撤去外力后平衡长度为 $l_p$。求 $l_p$ 的代数方程、平衡位置附近的等效劲度系数 $k_{\mathrm{eff}}$，以及可实现平衡的 $I_0$ 范围。

（2）改为零电阻理想导线闭合的回路，初始加载电流为 $I'_0$，撤去外力后平衡长度为 $l'_p$。求 $l'_p$ 与 $k'_{\mathrm{eff}}$。

\begin{figure}[htbp]
  \centering
  \begin{tikzpicture}[x=.68cm,y=.72cm,>=Stealth,font=\small]
  \begin{scope}
    \draw[thick] (-3.2,0) -- (-2.6,0);
    \draw[thick,decorate,decoration={coil,aspect=.38,segment length=4.3pt,amplitude=3pt}] (-2.6,0)--(2.1,0);
    \draw[thick] (2.1,0)--(2.9,0);
    \draw[->,thick] (-2.5,.62)--(2.0,.62) node[midway,above] {$I_0$};
    \draw[<->] (-2.58,-.55)--(2.08,-.55) node[midway,below] {$l$};
    \draw[->,thick] (-3.15,0)--(-4.0,0) node[left] {$F$};
    \draw[->,thick] (2.85,0)--(3.7,0) node[right] {$F$};
    \draw[thick] (-2.55,1.25)--(-2.55,2.1)--(-1.35,2.1);
    \draw[thick] (2.05,1.25)--(2.05,2.1)--(.85,2.1);
    \draw[thick] (-1.35,2.1) arc[start angle=180,end angle=0,radius=1.1];
    \draw[fill=white,thick] (-.75,1.72) rectangle (.25,2.48);
    \node at (-.25,2.1) {恒流源};
    \node[below] at (-.15,-1.35) {（a）恒流条件：$I=I_0$};
  \end{scope}
  \begin{scope}[xshift=10cm]
    \draw[thick] (-3.2,0) -- (-2.6,0);
    \draw[thick,decorate,decoration={coil,aspect=.38,segment length=4.3pt,amplitude=3pt}] (-2.6,0)--(2.1,0);
    \draw[thick] (2.1,0)--(2.9,0);
    \draw[<->] (-2.58,-.55)--(2.08,-.55) node[midway,below] {$l$};
    \draw[->,thick] (-3.15,0)--(-4.0,0) node[left] {$F$};
    \draw[->,thick] (2.85,0)--(3.7,0) node[right] {$F$};
    \draw[thick] (-2.55,1.25)--(-2.55,2.1)--(2.05,2.1)--(2.05,1.25);
    \draw[->,thick] (-1.8,2.1)--(.65,2.1) node[midway,above] {$I$};
    \node[above] at (-.25,2.12) {零电阻闭合回路};
    \node[below] at (-.15,-1.35) {（b）磁通守恒：$LI=L_0I'_0$};
  \end{scope}
\end{tikzpicture}

  \caption{第39届复赛第2题的解析辅助示意。原题与答案均无独立图示；左为恒流源维持电流，右为零电阻闭合回路。}
  \label{fig:cpho39q02}
\end{figure}

\subsection{参考解答整理}
\paragraph{（1）恒流条件。}
长度为 $l$、电流为 $I$ 时，螺线管的磁感应强度、磁通匝链数与电感分别为
\[
B=\frac{\mu_0NI}{l},\qquad
\Psi=\frac{\mu_0N^2\pi r^2I}{l},\qquad
L=\frac{\mu_0N^2\pi r^2}{l}.
\]
恒流源做功须计入等效势能。弹簧势能、磁能和恒流源项合并为
\[
W=\frac12k(l-l_0)^2-\frac{\mu_0N^2\pi r^2I_0^2}{2l}.
\]
故维持平衡所需外力为
\[
F=k(l-l_0)+\frac{\mu_0\pi N^2r^2I_0^2}{2l^2}.
\]
撤去外力后的平衡条件为
\[
\boxed{k(l_p-l_0)+\frac{\mu_0\pi N^2r^2I_0^2}{2l_p^2}=0.}
\]
令 $l=l_p+\delta l$，一阶展开得到
\[
\boxed{k_{\mathrm{eff}}=k-\frac{\mu_0\pi N^2r^2I_0^2}{l_p^3}.}
\]
由平衡方程和平均值不等式，恒流的范围为
\[
\boxed{0\le I_0\le\sqrt{\frac{8kl_0^3}{27\mu_0\pi N^2r^2}}.}
\]

\paragraph{（2）零电阻闭合回路。}
回路电阻为零时磁通守恒，满足
\[
LI=L_0I'_0=\frac{\mu_0\pi N^2r^2}{l_0}I'_0.
\]
此时磁能正比于 $l$，系统能量为
\[
W=\frac12k(l-l_0)^2+\frac{\mu_0N^2\pi r^2I_0'^2}{2l_0^2}l.
\]
平衡条件给出
\[
\boxed{l'_p=l_0-\frac{\mu_0\pi N^2r^2I_0'^2}{2kl_0^2},\qquad k'_{\mathrm{eff}}=k.}
\]

\subsection{转写依据}
题干依据 \texttt{CPHO39-Q02.md} 第2页；公式依据 \texttt{CPHO39-Q02\_formula\_check\_v1.csv} 与答案第6至8页。原题与答案页均无独立图示，图\ref{fig:cpho39q02}为依据题设电路条件绘制的解析辅助图。本题第（1）问26分，其中（i）21分、（ii）5分；第（2）问14分。
